\scanpage{138}
% archive sheet number: 137

dont la forme quadratique associée

\diagram{\begin{tikzcd}[ampersand replacement=\&,column sep=large]
(\underbrace{L\otimes L'}_{\substack{\text{esp. des}\\\text{translations}\\\text{de }\underline L_P}})^{-1}
\arrow[r,"{\deleted{\unsure{\delta_{P,B}^{0}|\delta_{P,B}}}}"] \& (L\otimes L')^{-2}
\end{tikzcd}}

est l'élévation au carré (donc $\delta_{P,B}$
est \deleted{qui pour la} un morphisme fini
surjectif de degré 2).

Prenons l'image inverse de la section
nulle, c'est un schéma quadratique
\deleted{\unsure{$\widetilde S$|fini}} sur $S=\mathrm{Spec}\,k$ (donc défini par Alg.
quadratique $k'$) \inserted{qui ne dépend que des \unsure{rev.|inv.|env.} quad.}, cela dit, on a
un \uline{morphisme canonique}

% slanted
\margin{$X,X'$ \unsure{de|des} $S$ définis par $H,H'$, \unsure{soient|soit|posé} $\widetilde S=X\wedge X'$}

\diagram{\fbox{\begin{tikzcd}[ampersand replacement=\&]
X\times_S X'\arrow[r,"\alpha"] \& \widetilde S
\end{tikzcd}}}

\deleted{\unsure{ce qui|tel que}} tel que

\[
\alpha(\bar x,x')=\alpha(x,\overline{x'})=\overline{\alpha(x,x')}
\]

(où $x\unsure{-|\mapsto}\bar x$ etc désigne les inv. canoniques)
d'où pour $X,X'$ étales (et alors $\widetilde S$ est étale)

\diagram{\begin{tikzcd}[ampersand replacement=\&]
\overbrace{X\wedge_S X'}^{\substack{\text{produit contracté}\\\text{de }(\mathbb Z/2\mathbb Z)_S\text{-torseurs}}}
\arrow[r] \& \widetilde S,
\end{tikzcd}}

\deleted{\unsure{défini}} hom de $(\mathbb Z/2\mathbb Z)_S$-torseurs, donc \uline{iso}.

Il faut définir l'application $\alpha$

Soient $\varphi:\deleted{\unsure{A}}H\unsure{-|\to} k,\quad \varphi':H'-k$

\deleted{deux hom. d'Alg, considérons la forme}

\[
\deleted{\varphi\otimes\varphi':H\otimes H'=P\longrightarrow k,\quad u\otimes v\mapsto\unsure{\varphi(u),\varphi(v)|\varphi(u)\varphi(v)|\varphi(u)\varphi'(v)}} 
\]

